\chapter{怎样改写基因}

\begin{kaichang}
先别碰实验器材，请你去厨房拿一杯酸奶。

酸奶是细菌做的：乳酸菌把牛奶里的乳糖变成乳酸，奶就变酸变稠。乳品厂里最吓人的事故，不是停电，也不是牛奶不够，而是一种比细菌还小几百倍的东西混进发酵罐——噬菌体，专门感染细菌的病毒。一旦它溜进去，几小时内就能把一整罐勤勤恳恳干活的乳酸菌杀得精光，奶还是奶，变不成酸奶，一家大工厂一天的损失可以用吨计算。

可是，乳品厂的工程师们很早就注意到一件怪事：有些乳酸菌株被某种噬菌体“欺负”过一次之后，就像记了仇一样，从此再也不怕这种病毒。细菌没有大脑，没有神经，没有记忆细胞——它拿什么“记仇”？

更让人想不到的是：这个谜团被解开之后，人类手里多了一件能改写几乎任何生物基因组的工具。你多半在新闻里听过它的名字：CRISPR。这一章，我们就从“细菌怎样记仇”出发，走到“人怎样改写基因”——并且要走到新闻标题通常不肯带你去的地方：这把剪刀究竟有多可靠。
\end{kaichang}

\section{从“读”到“写”：改写需要三件东西}

上一章（第 81 章）我们解决了“读”：用测序仪把 DNA 这本长书的字母一个个读出来。读一本书，只需要眼睛和耐心；可是要\textbf{改}一本书，你需要三样东西：

\begin{itemize}[itemsep=2pt]
\item \textbf{找到那一页}：在几十亿个字母里，准确定位你想改的那一处；
\item \textbf{动手改}：剪掉、换掉或插入字母；
\item \textbf{改完再读一遍}：确认改对了，也没碰坏别的地方。
\end{itemize}

这三件事，每一件都是独立的技术关卡。历史上，人类是先学会“用固定模具剪贴”，再学会“指哪剪哪”的。而这两代工具，居然都是跟细菌学的。

先说第一代。早在 CRISPR 之前几十年，科学家就发现细菌拥有另一套防病毒武器：\textbf{限制性内切酶}。噬菌体把自己的 DNA 注进细菌体内，这种酶就专门在特定的字母组合处下刀——比如有一种酶永远只认 GAATTC 这六个字母，见到就剪。病毒的 DNA 长书被剪成碎片，入侵就此失败。那细菌自己的 DNA 为什么没事？因为细菌给自己 DNA 上的同样位置贴了化学“护身符”（甲基化修饰，第 74 章讲过的表观遗传手段之一），酶认不出来。

1970 年代，科学家回过神来：既然这种酶永远在固定的字母组合处下刀，那它就是一把\textbf{认字的分子剪刀}。再给它配一瓶分子胶水——\textbf{连接酶}，能把两段 DNA 的断口重新接好。这里插一句能量账：断口接上意味着形成新的磷酸二酯键，按第 27 章的规矩，成键放能，但连接酶催化的整个反应仍要消耗 ATP 来推动（第 50 章），断键吸能、成键放能的老规矩在这里一条不少。

剪刀加胶水，人类就有了第一套“剪贴 DNA”的技术——\textbf{重组 DNA}：从一种生物的基因组里剪下一段基因，接进另一段 DNA 里。只不过，第一代剪刀的“认字表”是固定的：每种限制酶只认 4 到 8 个字母的固定图案，由不得你指定“剪在我想剪的地方”。所以在很长一段时间里，改写基因更像\textbf{用固定模具冲压零件}，而不是\textbf{拿笔在纸上写字}。

\section{质粒与转化：让细菌替我们抄写}

剪下来、接上去的 DNA 片段，还需要一个搬运工才能送进细胞、并且被大量复制。这个搬运工叫\textbf{质粒}：细菌体内一种环状的小 DNA，只有几千个字母，能独立于主染色体自我复制，像主书之外的一本可以自由传抄的小册子。

操作流程是这样的：在试管里用限制酶把质粒剪开，把目的基因用连接酶接进去，再让质粒“回到”细菌体内。最后这一步叫\textbf{转化}：用冷热刺激、短暂电击或化学处理，让细菌的细胞膜（第 51 章那层磷脂边界）暂时漏出小孔，质粒趁机溜进去。接下来，细菌每分裂一次，就把质粒也抄一份；如果质粒上带着一段蛋白质配方，细菌还会照着配方开工生产——一个活的“生物复印机加小工厂”。

问题来了：转化效率并不高，一皿细菌里只有一小撮真的收下了质粒，怎么把它们挑出来？工程师在质粒上预先装了一个\textbf{抗生素抗性基因}当“车票”：把细菌涂到含抗生素的培养基上，没有质粒的全部死掉，活下来的每个菌落都是“收了货”的。这正是第 76 章自然选择的逻辑——差异、选择压力、留强汰弱——只不过这一次，选择压力是人放的。

\begin{qiaoliang}
这套“让细菌抄书”的办法能把看不见的分子的放大到看得见吗？算一笔账。假设一种常用质粒在每个细菌里约有 500 个拷贝，培养液长到每毫升 $10^{9}$ 个细菌，取 1 升（$10^{3}$ 毫升）培养液，质粒分子总数大约是
\[500 \times 10^{9} \times 10^{3} = 5\times 10^{14}\ \text{个}.\]
用第 5 章的老朋友阿伏伽德罗常数 $\ud{6.02\times 10^{23}}{mol^{-1}}$ 换算，这大约是 $10^{-9}$ 摩尔——听起来微不足道。可一个约 5000 个碱基对的质粒，摩尔质量高达 $\ud{3\times 10^{6}}{g\,mol^{-1}}$ 上下，乘起来就是\textbf{几毫克}的 DNA——足够放在天平上称、放在电泳胶里看（第 81 章）。一个细菌小到你用显微镜才看得见，十亿个细菌加上它们的质粒，就能把“一个分子”的事业放大到“一毫克产品”的规模。这正是后来重组胰岛素等产业的出发点（具体故事在第 83 章）。
\end{qiaoliang}

\section{细菌的档案馆：CRISPR 本来是干什么的}

现在，回到酸奶厂的谜。细菌的“记仇本”，写在基因组里一段模样古怪的区域：一长串\textbf{重复—间隔—重复—间隔}的序列，像一堵砖墙。砖（重复序列）每块都一样，二三十个字母；砖与砖之间的灰缝（间隔序列）却各不相同。这个结构的名字很长：成簇的、规律间隔的短回文重复序列，缩写 CRISPR。

后来科学家把“灰缝”里的字母拿去跟病毒基因比对，答案让人起鸡皮疙瘩：\textbf{每一条灰缝，都是一小段病毒的 DNA}。原来，细菌的祖先遭遇噬菌体入侵时，会从敌人的基因组上剪下一小条“照片”，夹进这本档案馆（这一步叫\textbf{获得}）。等同一种病毒再次来犯，细菌就把这些间隔序列转录成一条条短 RNA——相当于印发的\textbf{通缉令}——交给一类叫 Cas 的蛋白质（比如大名鼎鼎的 Cas9）。Cas 蛋白握着通缉令在细胞里巡逻，一旦短 RNA 与入侵病毒的 DNA 互补配对成功，Cas9 立刻下刀，把病毒基因组拦腰剪断（这一步叫\textbf{干扰}）。

请注意这里的分子逻辑，和第 66 章的双螺旋、第 68 章的转录一脉相承：\textbf{一切定位都靠碱基互补配对}——A 找 T（RNA 里是 U）、G 找 C，靠氢键的数目和几何形状“认字”，没有任何魔法。Cas9 蛋白自己并不认字，认字的是那段大约 20 个字母的引导 RNA。

人类工程师看穿这一点之后，只做了\textbf{一步替换}：把细菌档案馆里抽出来的通缉令，换成我们自己设计的 20 个字母。Cas9 从此变成一把\textbf{指哪剪哪的可编程剪刀}。还有一个工程上极其重要的细节：Cas9 下刀处旁边必须挨着一小段叫 PAM 的标志序列（对最常用的 Cas9 来说是 NGG 三个字母），相当于剪刀要求的“停车位”。没有停车位，剪刀宁可不动——这也让误伤的机会降了一截，但只是降了一截，后面你会看到。

\section{剪断之后：修复方式决定结局}

很多新闻只讲“CRISPR 剪断了基因”，但切割只是开头。真正的“改写”，发生在细胞\textbf{修补断口}的时候。双链断裂对细胞是重大事故，必须立刻抢修，而细胞手里有两套主要修理工（第 67、72 章我们见过它们造成的后果）：

\begin{itemize}[itemsep=2pt]
\item \textbf{直接对接}（学名非同源末端连接）：把断口两头抓过来直接粘上。快，但毛糙——接的时候常常随手丢掉几个字母，或者多塞进去几个。对基因来说，字母数一变，后面的读码框全部错位（第 69 章），蛋白质指令变成乱码：这个基因被\textbf{敲除}了。所以只要剪一刀、等着细胞自己毛手毛脚地补，就有很大机会废掉目标基因。
\item \textbf{照模板修补}（学名同源定向修复）：如果断口附近存在一段和断口两侧同源的 DNA，细胞会把它当模板，精确地把断口修好。人为地同时送进一段“设计好的模板”，就能让细胞照葫芦画瓢，把特定字母换掉，甚至插入一整段新基因——这叫\textbf{敲入}。可惜，这套精修系统在多数细胞里效率不高，而且主要在细胞分裂期才开工，这正是“精确改写”至今仍然困难的原因之一。
\end{itemize}

\begin{center}
\begin{tikzpicture}
\draw[thick] (0,0) -- (4.4,0);
\draw[thick] (5.6,0) -- (10,0);
\draw[thick] (0,0.35) -- (4.4,0.35);
\draw[thick] (5.6,0.35) -- (10,0.35);
\node at (5,0.75) {\small 断口};
\draw[->,thick] (5,-0.2) -- (5,-1.1);
\node at (5,-1.5) {\small 细胞的两条抢修路线};
\draw[->,thick] (3.8,-1.9) -- (2.4,-2.7);
\node[align=center] at (2.2,-3.4) {\small 直接对接：丢/多几个字母\\ \small 读码框错位 $\rightarrow$ 基因敲除};
\draw[->,thick] (6.2,-1.9) -- (7.6,-2.7);
\node[align=center] at (7.8,-3.4) {\small 照模板修补：按设计换字母\\ \small 或插入片段 $\rightarrow$ 精确改写};
\end{tikzpicture}
\end{center}

上面这张图只是人为的表示法：两条横线代表 DNA 双链，箭头代表细胞的选择，真实的修复发生在纳米尺度、由好几种蛋白质协作完成，并没有什么岔路口指示牌。同一个断口走哪条路，取决于细胞类型、细胞周期和一点点概率——这也是为什么同一次编辑实验，不同细胞的结局可能不一样。

\begin{zhengju}
这套改写工具并不是某位天才在实验室里凭空发明的，它的故事从一句“不知道这是什么”开始。

1987 年，日本科学家石野良纯的团队研究大肠杆菌的一个基因时，注意到旁边有一串奇怪的重复序列：正读反读都一样的小段落，中间夹着五花八门的小片段。他们不知道这是什么，只能如实写进论文：“其生物学意义尚不清楚。”——科学的诚实，常常以这种不起眼的句子开头。

1990 年代，西班牙微生物学家莫希卡在极端嗜盐的古菌里也碰上了同样的结构。他追了十几年，发现这类序列在各种细菌和古菌里到处都是，而且间隔片段总在变化。2003 年他大胆推测：这可能和防御病毒有关。2005 年，三个研究组各自把大量间隔序列拿去比对数据库，发现许多间隔与噬菌体的 DNA \textbf{完全吻合}——“这是细菌的获得性免疫档案”的假说正式登场。但当时也有别的解释：也许只是巧合，也许这些序列另有调控功能。假说需要的不是更多人点头，而是一个能区分对错的实验。

关键实验在 2007 年到来，地点出人意料：不是名牌大学，而是丹麦一家乳品配料公司（Danisco）的实验室。巴朗古、霍瓦特等人用噬菌体攻击制作酸奶的嗜热链球菌，然后检查幸存者的基因组：幸存细菌的 CRISPR 区域里，多出了与攻击者精确匹配的新间隔序列。反过来，人为删掉一段间隔，细菌就丢掉对应的抗性；人为装进一段间隔，细菌就获得从没见过的抗性。干预一个变量，看到对应的因果变化——“档案馆记仇”的假说这才站稳。

机制层面的最后一块拼图在 2012 年补上：沙尔庞捷和杜德纳的团队把 Cas9 蛋白和两种 RNA 在试管里纯化出来重新组装，证明只需一条人工设计的引导 RNA，就能让 Cas9 在任意指定位置剪断 DNA——“可编程”三个字，在烧瓶里被直接演示。2013 年，张锋团队和丘奇团队几乎同时把这套系统搬进人和小鼠的细胞，成功改写了哺乳动物基因组。2020 年，沙尔庞捷和杜德纳因此获得诺贝尔化学奖。

从“没人认识的怪序列”到“改写生命的工具”，这条路走了 25 年：记录异常、比对数据库、提出假说、设计干预实验、在试管里重组机制。值得记住的不是谁最聪明，而是这条证据链的每一环，都曾经是可以被推翻的。
\end{zhengju}

\section{这把剪刀会剪歪：脱靶、嵌合与递送}

现在讲新闻标题常常省略的部分。CRISPR 被称作“基因剪刀”，这个比喻好用到容易骗人——剪刀这个词暗示着干净利落、指哪打哪。真实的编辑过程，至少要过三道鬼门关。

\textbf{第一关：脱靶。}引导 RNA 靠 20 个字母的互补配对认路，可“认路”是概率事件，不是非黑即白。人类基因组有 30 亿个字母，某段和通缉令只差一两个字母的“相似段落”，引导 RNA 有时也会凑合着配上对，Cas9 就在计划外的位置剪了一刀。这叫\textbf{脱靶}。脱靶的后果轻重不一：剪在不重要的区域也许平安无事，剪在另一个基因里就可能制造新的突变，极端情况下甚至激活癌症相关的通路（第 80 章讲过癌症是体内细胞的演化）。新一代 Cas 蛋白和改良的引导 RNA 都在拼命压低估靶率，但“压到低”不等于“降到零”。

\textbf{第二关：嵌合。}就算剪刀刀刀命中，也不是所有细胞同时被改。一次编辑处理的是千千万万个细胞：有的细胞成功改了，有的没改，有的改得五花八门——得到的组织或胚胎是一个“拼布被子”，不同格子里的字母不一样，这叫\textbf{嵌合}。在胚胎发育极早期做编辑，嵌合尤其常见：等修改完成时，胚胎已经分裂出好几个细胞，只有部分细胞带着想要的修改。所以“一个编辑过的胚胎”这个说法本身就很含糊——它可能只是“一部分细胞被编辑过的胚胎”。

\textbf{第三关：递送。}这是最不浪漫、却往往最难的一关：剪刀再好，得先\textbf{送进正确的细胞}。在体外，事情好办——把细胞从身体里取出来，用电击把编辑工具灌进去，改完再送回去（某些细胞疗法就是这个思路，第 83 章细讲）。可要在身体里直接改，就得给 Cas9 和引导 RNA 找“快递”：常用的是去掉毒性的病毒载体，或者脂质纳米颗粒（和第 47 章的脂质、第 51 章的膜是同一套化学）。肝脏血流丰富，快递相对好送；大脑有血脑屏障把守，极难送达；快递包装本身还可能触发免疫系统的警报（第 61 章），轻则失效，重则伤身。一项编辑技术“在试管里成功”和“在病人身上成功”，中间隔着的常常就是这一关。

\begin{wuqu}
“CRISPR 就像文字处理软件里的查找替换：输入想改的序列，就能精准改好。”

这个比喻错在三个地方。第一，查找替换\textbf{不会打错字}，而引导 RNA 容忍一两个字母的错配，可能改到相似段落上（脱靶）。第二，查找替换\textbf{不会只改一半文档}，而真实编辑的产物常常是嵌合体：一部分细胞改了，一部分没改。第三，你的文档\textbf{不需要快递}，而身体里的目标细胞要靠病毒载体或脂质颗粒才能触及。更贴切的画面是：派一个只认 20 个字、偶尔认错一两个字的快递员，进一座藏书 30 亿字、还在不停自我复印和自我修补的图书馆——成功率、准确率、送达率，是三场彼此独立的考试，每一场都可能挂科。把 CRISPR 说成“无误差基因剪刀”，是对这三场考试的全盘隐瞒。
\end{wuqu}

\begin{tiaozhan}
为什么 20 个字母的通缉令还会认错人？用第 81 章的组合直觉粗算一下。一段特定的 20 字母序列，随机出现在某处的概率约为 $4^{-20}\approx 10^{-12}$，乘上基因组的 $3\times 10^{9}$ 个字母，纯随机撞车几乎不可能——这正是“20 个字母足以定位”的理由。但脱靶发生在\textbf{近似匹配}上：只差 1 个字母的“邻居序列”有 $20\times 3=60$ 种写法，差 2 个字母的有 $\binom{20}{2}\times 3^2\approx 1700$ 种写法。把几千种可能的近似写法撒进 30 亿个字母里，碰上几个“七分像”的段落就不再是小概率怪事，而是可以预料的统计必然。这就是为什么工程师要花大力气改造 Cas 蛋白的“认字严格程度”——他们对抗的不是偶然事故，而是组合数学。这段推导跳过不影响主线，但它能让你对新闻里“脱靶率降低了多少倍”的报道有自己的判断力。
\end{tiaozhan}

\section{更精细的工具：碱基编辑与先导编辑}

剪断双链再赌细胞的修补手艺，终归有点粗暴。能不能\textbf{不剪断}，只换掉一个字母？2016 年以来出现的\textbf{碱基编辑}做到了这一点：工程师把 Cas9 的“刀刃”钝化——它还能被引导 RNA 带到指定位置，却不再剪断双链——然后在这个定位器上挂一个“改字母的酶”，把 C 改成 T，或者把 A 改成 G。不少遗传病的病根恰好就是\textbf{一个字母}之差，对这类病，碱基编辑像铅笔和橡皮，比剪刀温和得多，双链断裂带来的大片插入缺失几乎消失。它的短板也很明确：只能做特定几种字母转换，想改的字不在菜单上就没辙。

2019 年登场的\textbf{先导编辑}更进一步：Cas9 只轻轻剪开一条链，身上挂着逆转录酶，引导 RNA 加长并自带一小段“新文字模板”——机器就定位在原位，照着模板把新字母\textbf{抄}进基因组。如果说 CRISPR 加模板是“剪开书页、盼着誊写员按你给的样张重抄”，先导编辑就是“翻开那一页，当场把错字描掉重写”，理论上能做任意字母替换、小段插入和小段删除。但请注意它的年龄：效率还不够高，递送仍然困难，出错的可能依然存在——它只是把三道鬼门关的难度各调低了一档，并没有拆掉任何一道。

\begin{table}[htbp]
\centering
\caption{几代“改写基因”工具的对比（“像什么”只是帮助记忆的比喻）}
\begin{tabular}{llll}
\toprule
工具 & 像什么 & 擅长 & 主要短板 \\
\midrule
限制酶＋连接酶 & 固定模具＋胶水 & 拼接、克隆 DNA & 只能在固定位点下刀 \\
CRISPR 敲除 & 可编程剪刀 & 废掉目标基因 & 脱靶；修补结果杂乱 \\
CRISPR＋模板 & 剪刀＋誊写员 & 精确换字母、插入 & 效率低；依赖细胞分裂 \\
碱基编辑 & 铅笔与橡皮 & 单个字母转换 & 只能做几种转换 \\
先导编辑 & 搜索并替换 & 小范围任意改写 & 效率与递送仍受限 \\
\bottomrule
\end{tabular}
\end{table}

回头看，这张表的演进方向很清楚：定位越来越自由，动作越来越轻，留给细胞的“自由发挥空间”越来越小。但每一行“主要短板”列里，都还写着没写完的句子。

\section{边界：这套工具什么时候不好用}

任何工具都有适用边界，改写基因的工具尤其需要把丑话说在前面。

第一，它最擅长\textbf{因果清楚的单基因问题}：一个字母、一个基因、一种病，链条短，下刀处明确。可是身高、绝大部分常见病风险这类性状，是成百上千个基因各自贡献一点点、再和环境搅在一起的统计效应（第 73、85 章）——没有“那一刀”可剪，找遍基因组也找不到该改的地方。

第二，就算目标明确，\textbf{改体细胞}和\textbf{改生殖系}是两回事。改一个病人骨髓或肝脏里的细胞，修改随这个人离世而结束；改胚胎或生殖细胞，修改会写进往后所有世代。前者是治病，后者是替尚未出生、无法表态的人做决定——这条线怎么划，是技术问题更是伦理问题，第 83 章会专门讨论。

第三，改写字母不等于掌握了基因的\textbf{使用规则}。第 70 章讲过基因像可调光的灯而不是一直开着的灯，第 74 章讲过序列不变、使用方式也会变。把一个字母改对了，它在肝脏、在大脑、在十岁、在五十岁分别造成什么后果，常常是另一本还没读完的书。

\begin{shiyan}
\textbf{把 DNA 从水果里“请”出来}（课堂或在家都能做，全程不动基因，只是取出 DNA 看看它长什么样）。

材料：半根熟透的香蕉（或两颗草莓）、一小勺盐、两滴洗洁精、半杯温水、纱布或咖啡滤纸、一个透明杯子、一杯在冰箱里镇了半小时的医用酒精。

步骤：①香蕉加温水、盐、洗洁精，装进保鲜袋揉两三分钟成糊——盐让 DNA 更容易聚拢，洗洁精拆散细胞膜和核膜的磷脂（第 51 章）；②用纱布把糊过滤到杯子里；③把冰酒精沿杯壁缓缓倒进去，让它浮成一层，不要搅拌；④静置五分钟，盯着两层交界处。

观察点：交界处会慢慢出现白色絮状或丝状的东西，用筷子能绕起来——那就是几亿个香蕉细胞释放的 DNA 缠在一起的样子。

安全提示：酒精\textbf{易燃}，全程远离明火，用完盖紧；实验材料接触过洗洁精，\textbf{禁止入口}；做完用肥皂洗手。

想一想：你手里这团白色絮状物里就有香蕉的全部基因。可它离“改写一个基因”还差着定位、切割、修复、递送四道关——亲手摸一摸 DNA，再回头体会这条工具链有多长。
\end{shiyan}

\section{回到那杯酸奶}

现在可以回答开场的问题了。乳酸菌凭什么记仇？凭的是写在基因组里的档案馆：每熬过一次噬菌体入侵，就把敌人的一小段 DNA 夹进 CRISPR 区域；Cas 蛋白握着由这些片段转录出的通缉令日夜巡逻，仇人再度登门时当场剪碎。这是一套货真价实的获得性免疫系统——比脊椎动物的抗体早出现不知多少亿年，而且直接把“记忆”写进了可以遗传的 DNA。

而人类做的事情，从分子的角度看其实只有一件：\textbf{把通缉令换成自己写的}。细菌的防御档案，就这样变成了我们的编辑工具。顺带说，乳品工业今天仍在用这套机制本身：通过检查菌株的 CRISPR 档案，选育对流行噬菌体有抗性的发酵菌——你手里的这杯酸奶能稳定地产出，部分要归功于细菌几十亿年练出来的记仇本能。从酸奶到诺奖，从病毒碎片到改写生命的铅笔，这条路没有一个环节是凭空跳出来的。

\section{继续深挖：改写之后，问题才刚开始}

会改写基因，立刻引出两串更大的问题。第一串是医学的：这套工具怎样变成真正的药？从重组胰岛素到基因治疗，从“改体细胞治病”到“生殖系编辑”那条必须划清的线，再到把细胞取出来改好再送回去的 CAR-T 疗法——那是第 83 章的内容。第二串是生态与社会的：把改造过的生物种进田里、放进环境，基因会不会流到野外？“改写生命”的边界应该由谁来划？第 84 章会接着谈。你现在已经能看清一件事：真正稀缺的从来不是“剪一刀”的技术，而是知道\textbf{该不该剪、剪哪里、剪坏了怎么办}的知识与约束。

\section*{练一练}

\begin{sikaoti}
\item 画一幅信息流图：从“噬菌体首次入侵细菌”到“细菌获得针对它的抗性”，标出 DNA、RNA、蛋白质分别在哪一步登场，哪一步发生“获得”，哪一步发生“干扰”。
\item 定位为什么需要大约 20 个字母而不是 4 个？算一算：4 个字母的序列有 $4^4=256$ 种写法，在 30 亿字母的基因组里平均会出现多少次？20 个字母呢？（用 $4^{20}\approx 10^{12}$ 估算。）
\item 同样剪一刀，为什么大多数时候得到的是“基因敲除”，只有加了模板、运气又好时才得到“精确改写”？请从细胞两条修补水道的开工条件解释。
\item 有人宣称“我们用 CRISPR 改造了某种植物，绝对安全”。请设计一个核查思路，检验编辑是否发生在计划位置、有没有脱靶（提示：把这一章的工具和第 81 章的测序组合起来）。
\item 一个 4 细胞期的胚胎接受编辑，只有 1 个细胞被成功修改。粗略估计这个胚胎长成的个体中，多大比例的细胞携带修改？如果生殖细胞恰好不在其中，修改还会传给下一代吗？这说明“编辑了一个胚胎”这句话需要怎样更精确地表述？
\item 限制酶只认 GAATTC 这样的固定序列，细菌自己的基因组里也有这种序列。细菌靠什么保护自己？这种“给自家人贴标签”的思路，和第 61 章免疫系统区分“自己人”和“入侵者”有什么相似之处？
\end{sikaoti}
